% For Dutch text, use option `dutch'
% For VU students, use option `vu'
%\documentclass[12pt,vu]{adammath}
\documentclass[12pt,vu]{adammath}

% \usepackage[dutch]{babel}
\usepackage[english]{babel}
\usepackage{graphicx}
\usepackage[pdfborder={0 0 0}]{hyperref}
\usepackage{lipsum}

\usepackage{amsmath}
\usepackage{amssymb}  
\usepackage{amsthm}

\usepackage{color}
\definecolor{keywordcolor}{rgb}{0.7, 0.1, 0.1}   % red
\definecolor{commentcolor}{rgb}{0.4, 0.4, 0.4}   % grey
\definecolor{symbolcolor}{rgb}{0.0, 0.1, 0.6}    % blue
\definecolor{sortcolor}{rgb}{0.1, 0.5, 0.1}      % green

\usepackage{listings}
\def\lstlanguagefiles{lstlean.tex}
\lstset{language=lean}

%
% Data
%

\title{Logical Verification of a Compiler in LEAN}
\author[m.rimkevicius@student.vu.nl, 2708302]{Martynas Rimkevičius}
\what{Bachelor thesis}


\programme{Bachelor of Computer Science}
%\track{Dynamiische Systemen}

\supervisor{prof.dr.\ Jasmin Christian Blanchette}
\secondgrader{dr.\ xxxxx}
\coverimage{\includegraphics[scale=1.2]{lean}}

%
% Math environments
%

\theoremstyle{plain}                    
\newtheorem{theorem}{Theorem}[chapter]
\newtheorem{lemma}[theorem]{Lemma}           
\newtheorem*{maintheorem}{Proposition A}  % special theorem

\theoremstyle{definition}
\newtheorem{definition}[theorem]{Definition}  
\newtheorem{example}[theorem]{Example}

\theoremstyle{remark}
\newtheorem*{remark}{Remark}

%
% Maths
%

% examples
\newcommand{\field}[1]{\mathbb{#1}}
\newcommand{\R}{\field{R}} % real numbers
\newcommand{\N}{\field{N}} % natural numbers
\newcommand{\C}{\field{C}} % complex numbers
\newcommand{\Z}{\field{Z}} % integers

\newcommand{\ie}{\emph{i.e.,}\xspace}
\newcommand{\eg}{\emph{e.g.,}\xspace}
\newcommand{\etc}{etc.\xspace}
\newcommand{\etal}{\emph{et~al.}\xspace} 

\newcommand{\todo}[1]{\textcolor{blue}{#1}} 

%
% Content
%

\begin{document}
\maketitle

\begin{abstract}
% \lipsum[2-3]
\end{abstract}


\setcounter{secnumdepth}{3} % organisational level that receives a numbers
\setcounter{tocdepth}{3}    % print table of contents for level 3

\tableofcontents

\renewcommand{\chaptername}{} % uncomment to print only "1" not "Chapter 1"

\chapter{Introduction}\label{s:intro}



% Compilers are a crucial part of each computer and system since each static programming language relies on the correct abstraction of machine code and execution. This is where formal verification comes into effect. 

% Lean is a powerful tool to keep the proof correct since it verifies itself rather than a pen-and-paper proof. This means that mistakes are easy to spot and big-proof systems could be built on top of each other ground-up, which would not be as straightforward in a formal text-written proof. There are other theorem provers such as Isabelle/HOL, \todo{insert other provers}, but Lean has the advantage of \todo{insert why I am doing this in lean}

% In this paper, I will be creating a simple stack machine language - WHILE \cite{hithchiker}- compiler, and proving its correctness in LEAN. LEAN is a functional programming language that is used as a theorem prover to derive logical verification for its functions and types. \todo{expand}

% This thesis is heavily inspired by Chapter 8 of Concrete Semantics: With Isabelle/HOL \cite{isabelle}. It builds upon a simple imperative language defined in Chapter 8 of Hitchhiker's Guide to Logical Verification \cite{hithchiker}.

% \todo{talk about why and how I am contributing}


\chapter{Lean Proof Assistant}\label{s:lean}

In this section, I will go over the most relevant to the paper high-level Lean topics. The topics include type theory, inductive types and proof tactics. I will also compare them to Isabelle/HOL such that later proof comparisons would be more intuitive.

The Lean proof assistant is based on dependent type theory with inductive families.
%https://link.springer.com/article/10.1007/bf01211308 %

Simple type theory in Lean is used to associate every expression with a type. This means that there is no room for different interpretations. 
%https://leanprover.github.io/theorem_proving_in_lean/theorem_proving_in_lean.pdf
For example, each expression, a variable, a function or a constant, will have a type, for instance - an integer. What makes the Lean type theory powerful is that from the existing types we can build new ones. If $\alpha$ and $\beta$ are types, and $\alpha \rightarrow \beta$ is a function type out of them, there can also be a new type $\alpha \times \beta$.
%% Insert an example of that

Lean extends simple type theory by making treating the types themselves as objects. This means that types have to have a type which is the case - any type, even the basic such as \lstinline{nat} or \lstinline{bool}, have a type \lstinline{Type}.
%insert example from page 13
In proofs this allows us to operate on types the same way we operate on objects.

To define an object or a function of a certain type in Lean, \lstinline{def} is used to do that. Objects are used in proofs and proofs are also the ones that verify their correctness. The general form of a definition is as follows: $def foo : \alpha := bar$ (the type sometimes can be inferred, but the best practice is to specify it). 

Lean also has inductive types. Inductive types in Lean allow us to create types from scratch. A lot of definitions in Lean are made using inductive types, which allow us to use induction in the proofs or to define a recursive type. Inductive types are built from separate constructors, which are used to construct new objects of such type using previously defined values.
% (src: https://leanprover.github.io/theorem_proving_in_lean/inductive_types.html). 
The inductive type creates abstraction over the constructors and makes it easier to reason about each object collectively. This also allows to have recursion in the types themselves where the input becomes the type itself. This is because each constructor acts like a definition of that type, therefore it can have a its type passed as an argument.
% insert example of list

\begin{lstlisting}
    inductive list (T : Type) : Type :=
    | nil {} : list T
    | cons   : T → list T → list T
\end{lstlisting}

In Isabelle/HOL types are used in a similar manner. They are defined using the datatype keyword and they are built already as inductive types with possible several constructors.

\todo{add more text about this/add and explain examples}

Next important part is proof tactics in Lean. The tactics are recognized by keywords \lstinline{begin ... end} and \lstinline{by}. The tactic proofs are backward reasoning proofs, meaning the goal is broken down into smaller sections which are proven separately and construct a full proof. The advantage of these proofs are that the goal can be very complex and by splitting it into smaller problems, it makes it possible to achieve the goal. There are the opposite proof strategy - structured proofs, which are used for forward reasoning, are proofs which begin with hypotheses and assumptions and work towards the goal. In practice, both of these strategies are used together to split the goal into smaller parts and use hypotheses to prove subgoals.

Tactics in Lean can be partially automated, which are usually helpful to use to avoid long and tedious proofs. Some of the most used automated tactics in my case are \lstinline{simp} (it applies several trivial rewriting steps (can be by using a given hypothesis) which simplify the goal), \lstinline{linarith} (combines existing (in)equalities to find a contradiction between them), \lstinline{ring_nf} (solves equations that are in a commutative (semi)ring structure), \lstinline{finish} (simplifies the goal, eliminates quatifiers and looks for contradiction), and \lstinline{exact} (applies the hypothesis which is the same as the goal). Since most of these tactics help prove trivial goals or complete a set of trivial operations, they do not interfere with the clarity of the proof.

In Isabelle/HOL there are more automated proofs. There is \emph{simp} which is similar to Lean's, but there is also \emph{auto}, which is able to prove simple logical or set-theoretic goals \cite{isabelle}. Some of the automated proofs that are used more frequently are \emph{fastforce}, which tries harder than auto, and it works on the first subgoal only where it either succeeds or fails, and blast which is commonly used to prove complex logical goals. The automated proofs are said to be complete if they can prove all true formulas, but there is no complete proof method for Isabelle, which is why there are different automated proof methods.

\chapter{Stack machine instructions}\label{s:instr}

To set up a stack machine, we begin by defining the instruction set architecture and semantics. The definitions in this section are heavily inspired by Isabelle/HOL definitions in Concrete Semantics \cite{isabelle}, yet it is expanded upon and adapted to suit Lean's structure. Moreover, the stack machine is designed to fit a minimalist imperative language called WHILE defined in the Hitchhiker's Guide \cite{hithchiker}. 

We start by defining the basis of our stack machine - variables, values, and the stack. Since this is a simple stack machine, the supported value type is an integer and we can also store variables whose name is of string type. Stack in our case is created to save the values, and it is defined as a list of integers.
\begin{lstlisting}
    def vname := string
    def stack := list ℤ 
\end{lstlisting}

The way the variables are connected to the values is through a state function, which is essentially the machine's memory. State is a mapping of \lstinline{vname → ℤ} such that the value can be assigned and updated. I have used the state definition from Hitchhiker's Guide \cite{hithchiker}, where the state is defined as a mapping of \lstinline{string → ℕ} but since the variables are integers and the program counter is going to be integer, the definition was adapted accordingly. One of the key operations on the state that make it work as a mapping is the update operation.
\begin{lstlisting}
    def state.update (name : vname) (val : ℤ) (s : state) : state :=
    λname', if name' = name then val else s name'
\end{lstlisting}
This operation either assigns a new value to the existing variable or creates a new variable and does it each time without additional operations required. This is possible because the update has several proven lemmas, which are imported from the Hitchhiker's Guide \cite{hithchiker}, that can be used for further simplification in the application. 

The instruction set is comprised of simple, assembly-level instructions to put the value in the stack, do arithmetic operations on the values, change the variable, and change the order of execution. The definition of instruction is of previously-discussed inductive type, which makes each instruction a separate constructor.  
\begin{lstlisting}
    inductive instr : Type
    | LOADI : ℤ → instr
    | LOAD : vname → instr 
    | ADD : instr
    | SUB : instr
    | MUL : instr
    | DIV : instr
    | STORE : vname → instr
    | JMP : int → instr
    | JMPLESS : int → instr
    | JMPGE : int → instr
    | NOP : instr
\end{lstlisting}
The instruction \lstinline{LOADI} pushes the integer on the stack, \lstinline{LOAD} pushes the value of the variable on the stack, \lstinline{ADD}, \lstinline{SUB}, \lstinline{MUL}, \lstinline{DIV} do corresponding arithmetic operations with the top two integer values of the stack. \lstinline{STORE} instruction stores the top value of the stack in memory under a given variable name. \lstinline{JMP} jumps, or changes the program counter, by a given relative value. \lstinline{JPLESS} and \lstinline{JMPGE} compare the two top stack values and jumps if the second one is less or if the second one is greater or equal accordingly. 

In Isabelle the definitions of machine instructions and its architecture do not differ, because Isabelle does not differ as much when defining types and functions from Lean.
\chapter{Machine Executions}\label{s:exec}

Stack machine instructions are executed in steps, which are transitions between machine configurations. Machine configuration is defined as follows:
\begin{lstlisting}
    def config := ℤ × state × stack
\end{lstlisting}
Config is a tuple of a program counter, the current state or memory, and the current state of the stack.

The machine step, function \lstinline{iexec} is a transition from a given config to another one by a specific step. The function does the instruction and in the process either increases by one or changes the program counter to where the instruction says it to change to. For readability and later semantics proofs, I do not to create abbreviations for \lstinline{stk.tail.tail} since it represents itself well enough.
\begin{lstlisting}
def iexec : instr → config → config
| (LOADI n) (i, s, stk) := (i + 1, s, n :: stk)
| (LOAD v) (i, s, stk) := (i + 1, s, s v :: stk)
| ADD (i, s, stk) := (i + 1, s, (stk.tail.head + stk.head) :: stk.tail.tail)
| SUB (i, s, stk) := (i + 1, s, (stk.tail.head - stk.head) :: stk.tail.tail)
| MUL (i, s, stk) := (i + 1, s, (stk.tail.head * stk.head) :: stk.tail.tail)
| DIV (i, s, stk) := (i + 1, s, (stk.tail.head / stk.head) :: stk.tail.tail)
| (STORE v) (i, s, stk) := (i + 1, s{v ↦ (stk.head)}, stk.tail)
| (JMP n) (i, s, stk) := (i + 1 + n, s, stk)
| (JMPLESS n) (i, s, stk) := 
  if (stk.tail.head) < (stk.head)
  then (i + 1 + n, s, stk.tail.tail)
  else (i + 1, s, stk.tail.tail)
| (JMPGE n) (i, s, stk) := 
  if (stk.tail.head) ≥ (stk.head)
  then (i + 1 + n, s, stk.tail.tail)
  else (i + 1, s, stk.tail.tail)
| NOP (i, s, stk) := (i, s, stk)
\end{lstlisting}

Having this step defined is not enough to check whether it is a valid execution, since it just applies the instruction to the config. For that, I define a boolean check function, which checks whether, given a list of instructions (a program) and two configurations, there can be a single instruction step taken at the first configuration's program counter location. 
\todo{should i talk about isabelle syntax differences?}
\begin{lstlisting}
    def exec1  (li : list instr) (c : config)  (c' : config) : Prop := 
    ( c' = iexec (nth li c.fst) c ∧ 0 ≤ c.fst ∧ c.fst < li.length)  
\end{lstlisting}

\todo{would it be too talk about exec1I as an exec1 intro rule?}

At the same time, this function checks whether the first config's program counter is valid, or is in bounds of the instruction list. This check is important because the program counter is an integer, therefore could be less than zero, and it could also perform a wrong step if it is out of bounds. Since the program counter is an integer, I create a new definition of the nth item in the list, since in mathlib \cite{mathlib} there is only a definition of nth being a natural number. This is where the instruction \lstinline{NOP} comes in. It is defined such that getting an instruction would not be an option of the instr object, but the instr object itself, however, the option is required to define a case when the list is empty and thus returns option none. In my case, since it is very specific, the nth function works only on instruction lists.  
%https://leanprover-community.github.io/mathlib_docs/init/data/list/basic.html#list.nth}. 
\begin{lstlisting}
    def nth : list instr → ℤ → instr  
    | (a :: l) n := if (n = 0) 
      then a
      else nth l (n - 1)
    | list.nil n := NOP
\end{lstlisting}

After \lstinline{exec1} - checking whether there is a step between two configurations - the definition can be expanded to a check whether there exists a set of steps in a list of instructions between two configurations - \lstinline{exec}. This can be done by using reflexive transitive closure \lstinline{star} with the \lstinline{exec1} function. As the name suggests, \lstinline{star} is either a reflexive relation between, in this case, two configs or a transitive relation. Transitive relation means that there is a step or steps that connect the two configurations. 
\todo{insert lean code of exec?}

To continue building on top of the machine executions, they have to be functionally verified. This is done for singular and multiple-step executions, and mostly concerns expanding the instruction list forwards or backward. It is important to be able to append or prepend new instructions and know how it will change the execution because if a command of multiple instructions is created to be executed, and that command needs to be concatenated to another command, both of these commands' lists are concatenated.

I start by translating \lstinline{exec1_appendR} and \lstinline{exec_appendR} lemmas from Isabelle:
\begin{lstlisting}
    lemma exec1_appendR {li c c' li'} (h: exec1 li c c'): 
    exec1 (li ++ li') c c'
\end{lstlisting}
\begin{lstlisting}
    lemma exec_appendR {li c c' li'} (h: exec li c c'):
    exec (li ++ li') c c' :=
\end{lstlisting}

These lemmas show that if a new instruction list is appended to the right of the current list, the execution does not change because the program counter is pointing to the same initial list. To prove these lemmas, I create and prove lemma \lstinline{nth_append}, which applies \lstinline{nth} to the correct list according to the size of \emph{i}:
\begin{lstlisting}
    lemma nth_append {l1 l2 : list instr} {i : ℤ}
    (h_nneg : 0 ≤ i) :
    nth (l1 ++ l2) i = (
      if i < int.of_nat (list.length l1)
      then nth l1 i
      else nth l2 (i - list.length l1))
\end{lstlisting}
In Isabelle, the proof of \lstinline{nth_append} applies the induction on the first list length and uses \emph{auto simp}, which automatically proves this lemma. In Lean, however, the proof is based on induction but had to be manually split in cases to eliminate if-then-else structures and operate on existing inequalities. The proof is straightforward, since it covers simple cases, but needed some writing to get it done with the manual creation of hypotheses and rewrites.

Using \lstinline{nth_append} lemma the proof of \lstinline{exec1_appendR} is not easy but quite clear - show that the program counter is always in-bounds of the first list. To do so empty list case results in the false (from which any goal can be proven) since the program counter has to be in between 0 and the length of the list, which is 0, and the non-empty list case main idea is to use the hypothesis \emph{h} from which we know that according to \lstinline{exec1} definition: \lstinline{i < li.length}.

Lemma \lstinline{exec_appendR} is also based on induction (this time \lstinline{induction'} tactic, which generalizes induction hypothesis by default - useful in star proofs), just this time on star induction, which is the same as Isabelle, but I define the steps myself. In Isabelle, there is an automatic proof used - \emph{fastforce} - which is based on \emph{star.step} (reducing reflexive transitive closure to a singular step) and using \lstinline{exec1_appendR}. In Lean, using the singular \lstinline{star.single} (equivalent to Isabelle's \emph{star.step}) seems to be not enough, and I use \lstinline{star.tail} tactic, which splits the goal into a star step to a transitive config, and from that config a singular step to the end config.

The way there are two lemmas for appending the instruction list with a new list and preserving the execution order, there are two lemmas that show how the order is preserved when the lists are prepended (appended to the left). When the list is added to the left of the current list, the program counter has to be shifted by the length of the added list in the starting config and the ending config.
\begin{lstlisting}
lemma exec1_appendL {i i' :ℤ} {li li' s stk s' stk'}
(h_li : exec1 li (i, s, stk) (i', s', stk')) :
exec1 (li' ++ li) ((list.length li') + i, s, stk) ((list.length li') + i', s', stk')

lemma exec_appendL {i i' :ℤ} (li li' i s stk i' s' stk')
(h_single : exec li (i, s, stk) (i', s', stk')) :
exec (li' ++ li) (list.length li' + i, s, stk) (list.length li' + i', s', stk')
\end{lstlisting}
To prove \lstinline{exec1_appendL}, I use the \lstinline{nth_append} lemma and the definition of \lstinline{exec1}. It also requires to use \lstinline{by_cases} tactic, which is necessary to eliminate \lstinline{ite (i < 0)} that appears from \lstinline{nth_append}. Additionally, lemma \lstinline{iexec_shift} that is defined in Concrete Semantics is necessary to simplify the left append, and that lemma states that \lstinline{iexec} operation is not affected by the shift. 
\begin{lstlisting}
lemma iexec_shift {instr n i' s' stk' i s stk}:
((n+i',s',stk') = iexec instr (n+i, s, stk)) = ((i',s',stk') = iexec instr (i,s,stk))
\end{lstlisting}
This lemma is proven using induction on \lstinline{instr}, thus proving that this is true for all instructions. All cases could be simplified into two - shift when there is \emph{JMP} instruction, and shift when there is not. Both cases are proven the same way, except when there is a \emph{JMP} instruction the equality of program counters contains the size of the jump, so I use two supportive lemmas to reuse the same proofs. Both supportive lemmas state a similar thing - if the program counter changes a certain way, when shift is applied, it will change the same way, and they have similar mathematical simplification using a \lstinline{iff.intro} tactic since one side simplifies because of the other being true and vice-versa.

Out of the intuitive lemmas where the new instruction list is appended to the left or right of the execution, there are three more lemmas defined in Concrete Semantics. These lemmas are necessary for later use since they allow reasoning about \lstinline{exec} in a more simple manner than always expanding its definition to \lstinline{star}. 

The most used lemma in this project in later stages is \lstinline{exec_append_trans} lemma, which shows that the separate executions of two lists can become one execution, if the first execution has the same end config as the second execution has the start config.
\begin{lstlisting}
lemma exec_append_trans {li' li s stk  s' stk' s'' stk''} {i i' j'' i'': ℤ}
(h1: exec li (0, s, stk) (i', s', stk'))
(h2: int.of_nat (list.length li) <= i' )
(h3: exec li' (i' - list.length li, s', stk') (i'', s'', stk''))
(h4: j'' = list.length li + i''): 
exec (li ++ li') (0, s, stk) (j'', s'', stk'')
\end{lstlisting}
The lemma statement is a little bit packed, but it allows me to build transitive steps for two (or more) instruction lists concatenated together in later proofs, thus breaking down the goal into smaller subgoals. To prove this lemma I use \lstinline{star.trans} tactic, which simplifies easily using the initial hypotheses and previously define \lstinline{exec_appendR}. Isabelle proof is automatic, uses \emph{metis}, but it does specify that it uses \emph{star\_trans} rule.
\chapter{Expressions}\label{s:expressions}
I will define the compiler in three separate steps as it is done in the Concrete Semantics \cite{isabelle}, beginning with arithmetic expressions, then boolean expressions and then for commands. Proofs of the correctness of these expressions will follow in the next section.

\section{Arithmetic expressions}

Arithmetic expressions are similar to the instruction definitions because each expression corresponds to an instruction that is put in the list. In the defi
\begin{lstlisting}
def acomp : aexp → list instr
| (num n) := [LOADI n]
| (var x) := [LOAD x]
| (add e₁ e₂) := acomp e₁ ++ acomp e₂ ++ [ADD]
| (sub e₁ e₂) := acomp e₁ ++ acomp e₂ ++ [SUB]
| (mul e₁ e₂) := acomp e₁ ++ acomp e₂ ++ [MUL]
| (div e₁ e₂) := acomp e₁ ++ acomp e₂ ++ [DIV]
\end{lstlisting}

\section{Boolean expressions} \label{bexp}
Boolean expressions are a little bit trickier to define since it does not directly correspond to instructions. In this case boolean expressions compilation takes additional parameters compared to arithmetic, of which flag \lstinline{f} of type \lstinline{Prop} indicates what value of the boolean expression represents the true case, offset \lstinline{n} is used to specify how much does the program counter needs to be shifted by. 
\begin{lstlisting}
noncomputable def bcomp : bexp → Prop → ℤ → list instr 
| (Bc v) f n := if (v = f) then [JMP n] else []
| (Not b) f n :=  bcomp b (¬ f) n
| (And b1 b2) f n := 
    let cb2 := bcomp b2 f n,
      m := if (f = true) then int.of_nat (list.length cb2) else int.of_nat (list.length cb2) + n,
      cb1 := bcomp b1 false m
    in cb1 ++ cb2
| (Less a1 a2) f n := acomp a1 ++ acomp a2 ++ 
  (if (f = true) then [JMPLESS n] else [JMPGE n])
\end{lstlisting}

\lstinline{Bc} or boolean constant case is simple - if the constant is equal to the flag, then create a jump, otherwise nothing is returned. The \lstinline{Not} case is self-explanatory - it changes the flag to be the opposite, such that during the compilation it compares the negated case. \lstinline{Less} is creating either \lstinline{JMPLESS} or \lstinline{JUMPGE} instruction depending on the flag \lstinline{f}. The \lstinline{And} expression compilation builds a list of instructions, where if the first compilation is true, the second is executed, otherwise after the first compilation jump is set to be over the second and by n. There is no way to make these compilations parallel, so they have to be executed sequentially, which makes the definition of the expression complicated.  

When defining \lstinline{bcomp} cases I use \lstinline{noncomputable} flag, because I generate data that is asserting its existence by using \lstinline{let} to define variables in \lstinline{And} case. This is described in Theorem Proving in Lean \cite{theorem_proving_in_lean}.

\section{Commands}

The main idea of compiled commands is a program that will perform the same state transformation and end with a program counter at the \lstinline{(ccomp c).length}. During that transformation, commands can take multiple intermediate state steps or have intermediate stack values but in the end, will result in a compiled program.  

\begin{lstlisting}
noncomputable def ccomp : com → list instr
| com.SKIP := []
| (Assign x a) := acomp a ++ [STORE x]
| (Seq c1 c2) := ccomp c1 ++ ccomp c2
| (If b c1 c2) := (
  let cc1 := ccomp c1,
    cc2 := ccomp c2,
    cb := bcomp b false (list.length cc1 + 1)
  in cb ++ cc1 ++ (JMP (list.length cc2) :: cc2)
)
| (While b c) := (
  let cc := ccomp c,
    cb := bcomp b false (list.length cc + 1)
  in cb ++ cc ++ [JMP (-(list.length cb + list.length cc + 1))] 
)
\end{lstlisting}

The commands mostly are trivial except the ones that have branches - \lstinline{If} and \lstinline{While}. \lstinline{If} compiles to the instruction list, where if the boolean expression is true, \lstinline{cc1} (\emph{THEN} branch) follows the condition without jumps and afterward jumps over \lstinline{cc2}. Otherwise, \lstinline{cc2} (\emph{ELSE} branch) is executed by jumping from condition over the \lstinline{cc1} branch and the \lstinline{JMP} instruction. \lstinline{While} command is in a similar manner as \lstinline{If}, just in this case, instead of \emph{THEN} branch there is either exiting the loop if the condition is false, or jumping back to the condition after the while body.

The command compilation also uses the \lstinline{noncomputable} flag because of the same reason as \lstinline{bcomp} - using \lstinline{let} in the definition.



\chapter{Compilation correctness}\label{s:correctness}
To have a complete compiler, I prove for each expression the correctness statement, which essentially means that the compiled expression executes as intended. 


\section{Arithmetic expressions}
The correctness of arithmetic expressions (like all correctness proofs) follows an induction on the arithmetic expression. Induction in this case, inspired by Isabelle proof, generalizes stack such that it could be specified in each induction hypothesis as needed (however, I don't ever use this). These induction cases can be split into two parts - numbers and variables, and operations with them because the proofs for these grouped cases are very similar. Isabelle proof is based on induction but is completed automatically by \emph{fastforce+}.

\begin{lstlisting}
lemma acomp_correct {a : aexp}  {s : state}  {stk : stack} :
exec (acomp a) (0, s, stk) (list.length (acomp a), s, (eval a s) :: stk)
\end{lstlisting}

\subsection{Numbers and variables}

Since the proofs for numbers (constants) and variables are very similar, they can be covered together. The compilation creates a list of a single instruction - \lstinline{LOADI} or \lstinline{LOAD}. This means that a single step execution is enough to verify a single instruction, therefore I use \lstinline{star.single} lemma to simplify the goal, and I use \lstinline{exec1I} lemma to create \lstinline{exec1 [LOADI a] (0, s, stk) (1, s, a :: stk)} subgoal (in the \lstinline{var} case it is \lstinline{[LOAD a]}).

\subsection{Operations with arithmetic expressions}

To prove the correctness of arithmetic operations I used \lstinline{star.trans} tactic twice to show how three lists are concatenated together. Both \lstinline{acomp a1} and \lstinline{acomp a2} compilations come from the respective inductive hypotheses, and the only part that I manually had to create was the arithmetic operator instruction (e.g. \lstinline{[DIV]}) execution itself. To create arithmetic instruction execution I could use \lstinline{star.single} lemma to transform multiple steps into one \lstinline{exec1} step and use \lstinline{exec1I} to create the goal execution.

\section{Boolean expressions}
Boolean expressions compilation correctness is based on induction where it generalizes \lstinline{f} (boolean flag) and \lstinline{n} (offset used in a jump instruction), with initial hypotheses being that \lstinline{0 ≤  n}. Having arbitrary flag and offset helps in proofs where the induction hypothesis (later $IH$) is very simple for separate executions and I need to create a sequential process out of those hypotheses, or I need to specify \lstinline{f}. All of the case proofs involve splitting the goal using \lstinline{by_cases} tactic or \lstinline{cases classical.em}, which create a true and false case for a boolean expression, which, in turn, eliminates \lstinline{ite} (if-then-else) in the goal.

\begin{lstlisting}
lemma bcomp_correct (n: ℤ) ( b f s stk)
(h_nneg : 0 ≤  n) :
exec (bcomp b f n) (0, s, stk) (list.length (bcomp b f n) + (if (f = bval b s) then n else 0), s, stk)
\end{lstlisting}

\subsection{Not}
In Isabelle, the \emph{Not} case is not completely automatic, but it explicitly specifies the \emph{f} to be \emph{$\sim$f} (not f), and then \emph{fastforce} completes the proof.

In Lean, the proof is more involved. Inspired by the Isabelle proof, I first specify $IH$ to use \lstinline{¬f} instead of \lstinline{f}. Then I split the if into true case and false case to simplify the goal. In the true case, where \lstinline{b} evaluates to be equal to \lstinline{f}, simplifying the goal and $IH$ is enough to get the same result, and in the false case simplifying is enough with a manual proof that not equals translates into one side is equal to negated other side (\lstinline{¬f = bval b s → f = ¬ bval b s}).

\subsection{And}
In Isabelle, the \emph{And} case is more specified than other proofs, meaning that the splitting of cases and $IH$ specification is explicitly defined. After these mentions, \emph{fastforce} automatic steps finishes the case.

In Lean, as in \lstinline{Not} case, I took inspiration of what explicit information Isabelle proof gave me and used it in my favor. As described in the Concrete Semantics \cite{isabelle}, the cases split is on \lstinline{f = true} since the jump after first part of execution depends on the \lstinline{f} (shown in the definition \ref{bexp}). In the true case of this split $IH$ is specified to have \lstinline{f} as false and \lstinline{n} as \lstinline{int.of_nat(list.length (bcomp b2 f n))}, whereas in the false case \lstinline{n} is \lstinline{ int.of_nat (list.length (bcomp b2 f n)) + n} (\lstinline{int.of_nat} is used because the length of a list is a natural number). In both cases, however, the proof follows the same structure but differs in small part due to the \lstinline{f} being true or false. Since the execution is of two instruction lists built during the boolean compilation, I use \lstinline{exec_append_trans} lemma to split the goal into smaller subgoals and to show how these two lists combine into one. The first list part is relatively easy to prove using specified $IH$ and the second part requires more details to get to the goal. The second part has an if-then-else in its goal program counter, which means that another case split on the value of the first \lstinline{And} part is created. In one of the cases, the first part is false, which means that the conjunction is false, so it jumps over the second part, making it a reflexive execution. In the other case, the jump depends on the second part being true or false, which means that it needs to be executed to get its value, thus there is no jump. The proof of the \lstinline{And} case is complex because it has to specify all these boolean values and cover all combinations.

\subsection{Less}

The \lstinline{Less} proof is constructed out of two applications of \lstinline{exec_append_trans} lemma to add together two arithmetic expressions and the jump command. To create arithmetic expressions I could use already proven \lstinline{acomp_correct} lemma, which makes the proof process easier. To create jump instructions I manually create an execution with an if-then-else, such that I could simplify it and prove it in smaller cases rather than combining cases bottom-up. The \lstinline{ite} simplification cases are of similar structure where \lstinline{exec} is transformed into \lstinline{exec1} by \lstinline{star.single} lemma, and then relevant comparison simplification is applied which is straightforward. 

In Isabelle, this is one of the cases of boolean compilation that is proven completely automatically by \emph{fastforce}.

\section{Commands}
The command correctness proof shows the correct representation of the source program in the machine program. This means that in the machine code it should cause the same state change and it should terminate if and only if the source program terminates. Termination can be expressed as a program counter reaching the end of the compiled expression, i.e. \lstinline{i = (ccomp c).length}, if it started from 0. In the same way, stack should not change.

The state change can be expressed as \lstinline{(c, s) ⟹ t}, which is called big-step operational semantics. This can be read as "Starting with the state s, executing c will terminate with a state t". I use modified big-step semantics defined from Hitchhiker's Guide \cite{hithchiker} that is used to create a minimalist imperative programming language interpreter. 

In the case of the compiler, the semantics matter when we want to be sure that the source code simulates compiled code - then every final state created by the compiled code is represented in the source code. Thus, we know that all compiled code is trusted if it terminates. To know that it terminates we need to prove that the compiled code represents correctly the source code \cite{isabelle}. I will do only the latter part which is easier \todo{say that i have not done the small-step semantics part and why} and it can be understood as follows: "if the source program executes from state s to state t, then the compiled program will as well." 

\begin{lstlisting}
lemma ccomp_bigstep {c : com} { s : state } {t : state } (stk : stack) 
(h_step : (c, s) ⟹ t) :
exec (ccomp c) (0, s, stk) (list.length (ccomp c), t, stk)
\end{lstlisting}

The proof of this lemma is by induction on big-step semantics \lstinline{h_step}. I use \lstinline{induction'} to generalize induction hypotheses (easier to use them) and to get cases that represent are from big-step semantics (splits into smaller cases which is also easier to work with).

\subsection{Assign}
The \lstinline{Assign} command is constructed out of arithmetic expression and \lstinline{STORE} instruction, so  \lstinline{exec_append_trans} lemma is applied to split the two in separate cases. The arithmetic expression can be proven by applying \lstinline{acomp_correct} lemma and for \lstinline{STORE} instruction \lstinline{star.single} lemma is used. To create the instruction, a hypothesis that matches the goal pattern is manually defined and simplified such that is true. This proof is relatively simple compared to other command cases, and that can also be seen in Isabelle, where this case is proven by \emph{fastforce} and \emph{simp}.

\subsection{Sequence}
The command \lstinline{Seq} is compiled out of two commands. In this case, in the proof for each of the commands, there is a relevant $IH$. These induction hypotheses are applied to show both commands in the \lstinline{exec_append_trans} application. Since these execution hypotheses are generalized in a way where they follow the transitive config (because of \lstinline{induction'}), there is no further work to be done.

In Isabelle, the proof follows the same structure by showing transitivity from the first part's start config to the second part's end config from induction hypotheses. These parts are joined together by automatic proof \emph{blast} which uses \emph{star\_trans}, which is essentially the same as using \lstinline{exec_append_trans} in Lean. 

\subsection{If}

\lstinline{If} is split into two true and false cases by the \lstinline{induction'} which made the proof easier to create and manage. Both cases follow a similar structure of connecting the condition to \emph{then} case or condition to \emph{else} case. Condition is created using \lstinline{bcomp_correct} lemma, but the resulting cases had to be manually defined. In the true case, there is another application of \lstinline{exec_append_trans} lemma to create a \emph{then} command compilation joined with a jump over \emph{else} case. Jump is proven with \lstinline{star.single} lemma and simplification, where the \emph{then} case is created from the $IH$. On the other hand, the false case, because it has a starting program counter position after \emph{then} case and the jump instruction, can be constructed using \lstinline{exec_appendL_if} lemma. This lemma states that if the program counter is more than or equal to the list, that is appended to the left, length, the execution continues in the initial list, in this case, the initial list being jump instruction concatenated to \emph{else} case. To skip the jump instruction \lstinline{exec_cons_1} lemma can be used, which states that if a single instruction is added to the left of the list, the execution continues in the initial list bounds but everything is shifted by 1. 

In Isabelle, this proof is completely automatic using \emph{fastforce+}.


\subsection{While}

The \lstinline{While} case proof is the most involved proof of this project. In Lean, the \lstinline{induction'} separates it into while true and while false cases. 

While false case is simple and similar to if \emph{else} case, where the condition is concatenated with something that does is skipped past. However, in the while loop, if the condition is false, the execution stops, so using \lstinline{exec_append_trans} lemma the boolean expression, which is created using \lstinline{bcomp_correct}, and execution of nothing, which is created with \lstinline{star.refl}, are connected together.

While true case is complex because it is a loop, so there are three executions to be connected and the loop execution itself. Inspired by Isabelle proof of \lstinline{While} case, I have defined all executions separately and then used \lstinline{exec_append_trans} to connect them. All three execution steps can probably be proven together, but separating them using \lstinline{exec_appendR} or \lstinline{exec_appendL} makes it easier to understand the goal and work towards it. The three steps of execution are condition, then while do part, and jumping back to the condition after do part is done. All of the steps are relatively simple when simplified from the initial goal because the initial list is puzzling to look at. In each step, I define manual equalities that I use to rewrite the goal to contain \lstinline{+ 0} or some negative list length because otherwise the \lstinline{exec_appendL} lemma would not work (it requires something to be added to the program counter, not only appended list). After simplifying the goals in each execution, condition is finished using \lstinline{bcomp_correct} lemma, while do section is finished using $IH$ of the singular execution, and jump back is finished with \lstinline{star.single} application. To complete the while, there is a the last part which is the loop execution itself, which is the compilation of the while. It might seem confusing why there is this execution of the while, but after a singular loop execution, the loop continues, thus it follows the singular execution. To prove the top-level while that follows individual execution steps, while $IH$ can be applied.
\chapter{Related Work}\label{s:related}

In this section, I will cover similar formal verification systems to what I have covered in this paper. This includes verification in proof assistants such as Isabelle \cite{isabelle, webasm}, Coq \cite{coq_concrete_semantics, coq_c}, program verifiers Why3 \cite{why3} and Stainless \cite{stainless}.  

Automatic proof assistants include Isabelle and Coq, in which the proofs are automated to the extent that the user does not need to manually create full proof systems. One of the works comparing these theorem provers is by Czajka \emph{et al.} \cite{coq_concrete_semantics}. They compare the assistants by translating theorems and proofs from Isabelle to Coq from Concrete Semantics by Nipkow and Klein \cite{isabelle} like I do to Lean but do that with additional chapters, which in my case have already been covered in Hitchhiker's Guide by Baanen \emph{et al.} \cite{hithchiker}. Their goal is to improve the CoqHammer tool, which searches for Coq proofs while showing the similarities and differences to Isabelle. Similar to Isabelle's auto or other automated tactics, CoqHammer has the main tactic \emph{hammer}, which combines machine learning and automated reasoning techniques to accomplish goals automatically. It also provides \emph{sauto} which is a strong version of Coq tactic \emph{auto} to simplify the goal as much as possible and perform rewrites using existing hints from an importable database. Even though Coq has automation, Czajka \emph{et al.} \cite{coq_concrete_semantics} show that Isabelle does perform faster (on the same hardware) and has fewer lines of code and tactics used than Coq. 

I have shown a small toy-like example of formal verification, but there exist real-world examples of formally verified systems. One of the examples is C11 compiler formalization in Coq by Krebbers \cite{coq_c}. The compiler is formalized in normalized CH$_2$O semantics, which is an incomplete but sufficient representation of the compiler in Coq. The normalized semantics were created such that if there is proof for any given program with respect to them, the proof would hold for any existing C compiler that adheres to the standard. It also allows improvement in the future versions of the C compiler standard and addresses issues on indeterminate memory and pointers which is important in such language as C. Another example of a real-world application of this is WebAssembly mechanization and verification in Isabelle by Watt \cite{webasm}. The verification of this programming language allowed Watt to fix several errors in the official WebAssembly specification and create an executable verified interpreter that could be used instead of the official interpreter in the future.

In addition to proof assistants, there are program verifiers that create dedicated environments to verify the behavior of programs written in some specific programming language. They can help one to verify programs written in some imperative language, and these program verifiers have a high level of automation in proofs because the languages specified in them are less expensive then in proof assistants. One of the program verifiers Why3 is discussed by Clochard \emph{et al.} \cite{why3}. Why3 was also used to create an interpreter and prove the correctness of a simple programming language, and it is designed to be even more automatized and fast. Automatization is relevant when proving not a small compiler but a complex program. Another program verifier example is Stainless which is a Scala program verifier developed by the LARA group in EPFL presented by Voirol \cite{stainless}. To verify them Scala programs are translated into an intermediate verification language using an Inox interpreter and then programs can be verified with some automation in the system using Stainless.


\chapter{Conclusion}\label{s:conclusion}


% \todo{
% Briefly summarize your contributions, and share a glimpse of the implications of
% this work for future research.
% }


% \lipsum[1-2]



%
% Bibliography
%

% The argument of the thebibliography environment (in this example `x') 
% denotes the length of the longest label. 
% If you have fewer than 10 references the argument is one character; 
% if you have 10 - 99 references you take an argument of two characters.
% \begin{thebibliography}{x}
% % Put bibliography in table of contents
% \addcontentsline{toc}{chapter}{Bibliography} 

% \bibitem{mascot}
%   H R Simpson and K Jackson,
%   \emph{Process Synchronisation in MASCOT},
%   The Computer Journal, \textbf{22}, 332,
%   1979.

% \bibitem{piethazegras}
%   \href{http://www.pythagoras.nu}{http://www.pythagoras.nu}
% \end{thebibliography}
\newpage
\addcontentsline{toc}{chapter}{Bibliography} 
\bibliographystyle{plain}
\bibliography{bib}

%
% Appendix
%

\appendix

% \chapter{Experiments}
% \lipsum[2]

\end{document}